\documentclass[10pt]{article}
\usepackage[T1]{fontenc}
\usepackage[utf8]{inputenc}
\usepackage[french]{babel}
\usepackage{amsmath,amssymb,mathrsfs,mathpazo}
\usepackage{geometry}
\geometry{a4paper,margin=2cm}
\usepackage{graphicx,array,tabularx,tabularray,colortbl,multirow,multicol,booktabs,enumitem,pifont,ulem,fancybox,diagbox,cancel}
\usepackage{pstricks,pstricks-add,pst-plot,pst-tree,pst-node,pst-text,pst-func,pst-math,pst-bspline,pst-3dplot,pst-eucl,pst-solides3d}
\let\pscalculate\undefined
\usepackage{pst-bezier}
\usepackage{tikz,pgfplots,tkz-tab}
\usetikzlibrary{arrows,arrows.meta,patterns,calc,positioning,trees,decorations.markings,shapes.geometric,shapes.misc,shapes.arrows,intersections,decorations.pathmorphing,decorations.pathreplacing,backgrounds,shadows,babel}
\pgfplotsset{compat=1.18}
\usepgfplotslibrary{fillbetween}
\usepackage[np]{numprint}
\usepackage{listings,makecell,pas-tableur,twoopt}
\let\barre\undefined
\usepackage[pstricks]{tablvar}
\definecolor{darkgreen}{RGB}{0,100,0}
\providecommand{\red}{\color{red}}
\providecommand{\blue}{\color{blue}}
\providecommand{\green}{\color{green}}
\providecommand{\black}{\color{black}}
\providecommand{\white}{\color{white}}
\providecommand{\gray}{\color{gray}}
\providecommand{\cyan}{\color{cyan}}
\providecommand{\magenta}{\color{magenta}}
\providecommand{\yellow}{\color{yellow}}
\providecommand{\violet}{\color{violet}}
\providecommand{\orange}{\color{orange}}
\providecommand{\euro}{\texteuro}
\providecommand{\EUR}{\texteuro}
\providecommand{\ieme}{\textsuperscript{e}}
\providecommand{\degres}{\ensuremath{^\circ}}
\providecommand{\degree}{\ensuremath{^\circ}}
\providecommand{\vv}[1]{\overrightarrow{#1}}
\providecommand{\Coord}[1]{#1}
\providecommand{\decosix}{}
\providecommand{\starredbullet}{\ensuremath{\bullet}}
\pagestyle{empty}
\setlength{\parindent}{0pt}
\usepackage{adjustbox}\DeclareUnicodeCharacter{2236}{:}\DeclareUnicodeCharacter{2190}{\ensuremath{\leftarrow}}
\providecommand{\R}{\mathbb{R}}
\providecommand{\N}{\mathbb{N}}
\providecommand{\Z}{\mathbb{Z}}
\providecommand{\Q}{\mathbb{Q}}
\providecommand{\D}{\mathbb{D}}
\providecommand{\C}{\mathbb{C}}
\providecommand{\e}{{\mathrm{e}}}
\providecommand{\vect}[1]{\overrightarrow{#1}}
\providecommand{\vectt}[1]{\overrightarrow{\mathrm{#1}}}
\providecommand{\vvt}[1]{\overrightarrow{\mathrm{#1}}}
\providecommand{\np}[2][]{{#2~#1}}
\providecommand{\vv}[1]{{\overrightarrow{#1}}}
\renewcommand{\d}{{\mathrm{d}}}
\providecommand{\degres}{{^\circ}}
\providecommand{\degre}{{^\circ}}
\providecommand{\blue}{\color{blue}}
\providecommand{\red}{\color{red}}
\providecommand{\enskip}{\quad}
\providecommand{\llbracket}{\lbrack\!\lbrack}
\providecommand{\rrbracket}{\rbrack\!\rbrack}
\providecommand{\euro}{{\text{€}}}
\providecommand{\no}{{\text{n°}}}
\providecommand{\interval}[2]{{[#1;#2]}}
\providecommand{\intervFO}[2]{{[#1;#2[}}
\providecommand{\dotfill}{\ldots}

\providecommand{\ds}{\displaystyle}
\providecommand{\textding}[1]{#1}
\providecommand{\Oij}{\(\mathrm{O};\vec{\imath},\vec{\jmath}\)}
\providecommand{\Oijk}{\(\mathrm{O};\vec{\imath},\vec{\jmath},\vec{k}\)}
\providecommand{\Ouv}{\(\mathrm{O};\vec{u},\vec{v}\)}
\providecommand{\intervFF}[2]{\ensuremath{[#1;#2]}}
\renewcommand{\euro}{\ifmmode\text{\texteuro}\else\texteuro\fi}
\DeclareUnicodeCharacter{00B0}{\ifmmode^{\circ}\else\textdegree\fi}
\DeclareUnicodeCharacter{20AC}{\ifmmode\text{\texteuro}\else\texteuro\fi}
% Énoncés : Éducation nationale. Transcription et corrigés : APMEP, Denis Vergès et contributeurs. Découpage et classement : CoopMaths, Sébastien Lozano, Éric Elter et contributeurs. Adaptation de la présentation : Maths Explorer.
% CC BY-SA : https://coopmaths.fr/www/about/
% https://coopmaths.fr/alea/static/bac/2023/tex/bac_2023_03_sujet1_reunion_4_cor.tex
\begin{document}
\noindent\textit{APMEP / CoopMaths. CC BY-SA. Présentation adaptée par Maths Explorer.}\par\medskip
\medskip

On considère le cube ABCDEFGH ci-dessous tel que AB $= 1$.

On note M le centre de la face BCGF et N le centre de la face EFGH.

\begin{center}
\psset{unit=1cm}
\begin{pspicture}(6,5.6)
\psframe(0.3,0.3)(4,4)
\psline(4,0.3)(5.7,1.4)(5.7,5.1)(4,4)
\psline(5.7,5.1)(2,5.1)(0.3,4)
\psline[linestyle=dashed](0.3,0.3)(2,1.4)(2,5.1)
\psline[linestyle=dashed](2,1.4)(5.7,1.4)
\uput[u](2,5.1){A} \uput[u](5.7,5.1){B} \uput[r](5.7,1.4){C} \uput[dr](2,1.4){D}
\uput[l](0.3,4){E} \uput[u](4,4){F} \uput[dr](4,0.3){G} \uput[dl](0.3,0.3){H}
\end{pspicture}
\end{center}

\medskip

On se place dans le repère orthonormé $\left(\text{D}~;~ \vect{\text{DH}},~\vect{\text{DC}},~\vect{\text{DA}}\right)$.

\medskip

\begin{enumerate}
\item 
On a: F\;$\begin{pmatrix} 1 \\ 1\\ 1\end{pmatrix}$ et C\;$\begin{pmatrix} 0 \\ 1\\ 0\end{pmatrix}$

\item 
Le point M est le centre de la face BCGF donc c'est le milieu de [CF]; il a donc pour coordonnées la demi-somme des coordonnées de C et de F:

M a pour coordonnées 
$\begin{pmatrix} 
\frac{x_{\text C} + x_{\text F}}{2} \\  
\frac{y_{\text C} + y_{\text F}}{2}\\  
\frac{z_{\text C} + z_{\text F}}{2}
\end{pmatrix}
=
\begin{pmatrix} 
\frac{0+1}{2} \\  
\frac{1+1}{2}\\  
\frac{0+1}{2}
\end{pmatrix}
=
\begin{pmatrix} 
\frac{1}{2} \\  
1\\  
\frac{1}{2}
\end{pmatrix}$

Le point N est le centre de la face EFGH donc c'est le milieu de [FH]; il a donc pour coordonnées la demi-somme des coordonnées de F et de H:

N a pour coordonnées 
$\begin{pmatrix} 
\frac{x_{\text F} + x_{\text H}}{2} \\  
\frac{y_{\text F} + y_{\text H}}{2}\\  
\frac{z_{\text F} + z_{\text H}}{2}
\end{pmatrix}
=
\begin{pmatrix} 
\frac{1+1}{2} \\  
\frac{1+0}{2}\\  
\frac{1+0}{2}
\end{pmatrix}
=
\begin{pmatrix} 
1 \\  
\frac{1}{2}\\  
\frac{1}{2}
\end{pmatrix}$

\item 
	\begin{enumerate}
		\item 
Le plan (HFC) a pour vecteurs directeurs $\vectt{HF}$ et $\vectt{CF}$ car les points H, F et C ne sont pas alignés.		

$\vectt{HF}$ a pour coordonnées
$\begin{pmatrix}
1-1\\1-0\\1-0
\end{pmatrix}
=
\begin{pmatrix}
0\\1\\1
\end{pmatrix}$
et
$\vectt{CF}$ a pour coordonnées
$\begin{pmatrix}
1-0\\1-1\\1-0
\end{pmatrix}
=
\begin{pmatrix}
1\\0\\1
\end{pmatrix}$
		
$\vectt{AG}$ a pour coordonnées
$\begin{pmatrix}
1-0\\1-0\\0-1
\end{pmatrix}
=
\begin{pmatrix}
1\\1\\-1
\end{pmatrix}$		
	
$\vectt{AG}\cdot\vectt{HF}=1\times 0 + 1\times 1 + (-1)\times 1 = 0$
donc $\vectt{AG}\perp\vectt{HF}$
		
$\vectt{AG}\cdot\vectt{CF}=1\times 1 + 1\times 0 + (-1)\times 1 = 0$
donc $\vectt{AG}\perp\vectt{CF}$		

Le vecteur $\vectt{AG}$ est orthogonal à deux vecteurs directeurs du plan (HFC) donc c'est un vecteur normal à ce plan.
		
		\item 
Le plan (HFC) passe par le point F et a pour vecteur normal $\vectt{AG}$, c'est donc l'ensemble des points X de coordonnées $(x\;,\;y\;,\;z)$ tels que 
$\vectt{FX}\perp\vectt{AG}$.		
		
$\vectt{FX}$ a pour coordonnées
$\begin{pmatrix}
x-1\\y-1\\z-1
\end{pmatrix}$
et $\vectt{AG}$ a pour coordonnées
$\begin{pmatrix}
1\\1\\-1
\end{pmatrix}$

$\begin{gathered}\vectt{FX}\perp\vectt{AG}
\iff \vectt{FX}\cdot\vectt{AG}=0
\iff (x-1)\times 1 + (y-1)\times 1 + (z-1)\times(-1)=0\\
\phantom{\vectt{FX}\perp\vectt{AG}}
\iff x-1+y-1-z+1=0
\iff x+y-z-1=0\end{gathered}$

Le plan (HFC) a donc pour équation cartésienne: $x+y-z-1=0$.

	\end{enumerate}
\item La droite (AG) a pour vecteur directeur $\vectt{AG}\;\left (1\;,\;1\;,\;-1 \right )$ et passe par le point A de coordonnées $\left ( 0\;,\;0\;,\;1\right )$; elle a donc pour représentation paramétrique:

$\left \{
\begin{array}{l !{=} r l}
x & x_{\text A} + 1\times t\\
y & y_{\text A} + 1 \times t & t\in\R\\
z & z_{\text A} + (-1)\times t
\end{array}
\right .$
soit
$\left \{
\begin{array}{l !{=} r l}
x &  t\\
y &  t & t\in\R\\
z & 1- t
\end{array}
\right .$

\item 
Soit R le projeté orthogonal de G sur le plan (HFC). donc R appartient au plan (HFC).

De plus, comme la droite (AG) est orthogonale au plan (HFC), le point R appartient à la droite (AG).

Les coordonnées de R vérifient donc le système:
$\left \{
\begin{array}{l !{=} r}
x &  t\\
y &  t \\
z & 1- t\\
x+y-z-1 & 0
\end{array}
\right .$

On a donc $t+t-(1-t)-1=0$ donc $3t=2$ donc $t=\dfrac{2}{3}$.

Donc R a pour coordonnées $x_{\text R}=t=\dfrac{2}{3}$, $y_{\text R}=t=\dfrac{2}{3}$, $z_{\text R}=1-t=\dfrac{1}{3}$, soit $\left (\dfrac{2}{3}~,~\dfrac{2}{3}~,~\dfrac{1}{3}\right )$.
\item Une représentation paramétrique de la droite (FG) est:
$\left\{\begin{array}{l !{=} l l}
x&1\\
y&1 & t\in\R\\ 
z&t
\end{array}\right.$.

On cherche un point K sur la droite (FG) tel que le triangle KMN soit rectangle en K.

K est un point de (FG) donc ses coordonnées sont de la forme $\left (1~,~1~,~t\right )$.

De plus, le triangle KMN est rectangle en K donc $\vectt{MK}\perp\vectt{NK}$ donc $\vectt{MK}\cdot\vectt{NK}=0$.

$\vectt{MK}$ a pour coordonnées
$\begin{pmatrix}
1-\frac{1}{2}\\1-1\\t-\frac{1}{2}
\end{pmatrix}
=
\begin{pmatrix}
\frac{1}{2}\\0\\t-\frac{1}{2}
\end{pmatrix}$
et
$\vectt{NK}$ a pour coordonnées
$\begin{pmatrix}
1-1\\1-\frac{1}{2}\\t-\frac{1}{2}
\end{pmatrix}
=
\begin{pmatrix}
0\\\frac{1}{2}\\t-\frac{1}{2}
\end{pmatrix}$

$\vectt{MK}\cdot\vectt{NK}=0
\iff \dfrac{1}{2}\times 0 + 0\times \dfrac{1}{2} + \left (t-\dfrac{1}{2}\right )\times \left (t-\dfrac{1}{2}\right ) = 0
\iff \left (t-\dfrac{1}{2}\right )^2=0
\iff t=\dfrac{1}{2}$

Le point K a donc pour coordonnées $\left (1~,~1~,~\dfrac{1}{2} \right )$.

\item 
Le tétraèdre FNKM a pour hauteur [FK] et pour base le triangle KMN.

\begin{list}{\textbullet}{}
\item F a pour coordonnées $\left (1~,~1~,~1\right )$ et K a pour coordonnées $\left (1~,~1~,~\dfrac{1}{2} \right )$; donc $\text{FK}=\dfrac{1}{2}$.

\item L'aire du triangle rectangle KMN est $\dfrac{\text{KM}\times\text{KN}}{2}$.

\begin{list}{$\circ$}{}
\item M\;$\left (\dfrac{1}{2}~,~1~,~\dfrac{1}{2} \right )$ et K\;$\left (1~,~1~,~\dfrac{1}{2} \right )$ donc $\text{KM}=\dfrac{1}{2}$
\item N\;$\left (1~,~\dfrac{1}{2}~,~\dfrac{1}{2} \right )$ et K\;$\left (1~,~1~,~\dfrac{1}{2} \right )$ donc $\text{KN}=\dfrac{1}{2}$
\end{list}

L'aire du triangle KMN est donc $\dfrac{\frac{1}{2}\times\frac{1}{2}}{2}=\dfrac{1}{8}$.

\item Le volume d'un tétraèdre  est, en unité de volume: 
$\dfrac{\text{aire de la base}\times\text{hauteur}}{3}$.

Le volume du tétraèdre FNKM est donc:
$\dfrac{\left (\text{aire de KMN}\right )\times\text{FK}}{3}
= \dfrac{\frac{1}{8}\times\frac{1}{2}}{3} = \dfrac{1}{48}$.
\end{list}

Le volume du cube est 1, celle du tétraèdre FNKM est $\dfrac{1}{48}$ donc le volume du tétraèdre représente un quarante-huitième du volume du cube.

\end{enumerate}

\begin{center}
\psset{unit=1cm}
\begin{pspicture}(6,5.6)
\psframe(0.3,0.3)(4,4)
\psline(4,0.3)(5.7,1.4)(5.7,5.1)(4,4)
\psline(5.7,5.1)(2,5.1)(0.3,4)
\psline[linestyle=dashed](0.3,0.3)(2,1.4)(2,5.1)
\psline[linestyle=dashed](2,1.4)(5.7,1.4)
\uput[u](2,5.1){A} \uput[u](5.7,5.1){B} \uput[r](5.7,1.4){C} \uput[dr](2,1.4){D}
\uput[l](0.3,4){E} \uput[u](4,4){F} \uput[dr](4,0.3){G} \uput[dl](0.3,0.3){H}

\psset{linecolor=blue,linewidth=1.3pt}
\psdots(4.85,2.7)(4,2.15)(2.15,2.15) 
\uput[dr](4.85,2.7){\blue{}M} \uput[dr](4,2.15){\blue{}K} \uput[l](2.15,2.15){\blue{}N}
\psline(4,4)(4.85,2.7)(4,2.15)(4,4)(2.15,2.15)(4,2.15)
\psline[linestyle=dashed](4.85,2.7)(2.15,2.15)
\end{pspicture}
\end{center}

\end{document}
